\documentclass[aspectratio=169]{beamer}
\usetheme{metropolis}
\usepackage{appendixnumberbeamer}
\usepackage{booktabs, hyperref, graphicx}

\input{mgmt675-style}

\usepackage{tikz}
\usetikzlibrary{shapes.geometric, arrows.meta, positioning, calc}

% Title info
\subtitle{MGMT 675: Generative AI for Finance}
\title{Module 8: Model Context Protocol}
\author{Kerry Back}
\date{}

\begin{document}

\maketitle

\begin{frame}{What is MCP?}
The \textbf{Model Context Protocol (MCP)} is an open standard that lets AI applications connect to external tools through a uniform interface.
\vspace{0.3cm}
\begin{columns}[T]
\begin{column}{0.5\textwidth}
\begin{shadedbox}[title=\textbf{The Idea}]
\begin{itemize}\small
  \item A standard ``plug-in'' protocol for AI
  \item Each server advertises tools (functions the AI can call)
  \item Tools appear automatically --- no custom code needed
\end{itemize}
\end{shadedbox}
\end{column}
\begin{column}{0.5\textwidth}
\begin{shadedbox}[title=\textbf{Analogy}]
\begin{itemize}\small
  \item USB is a standard for peripherals
  \item Plug in a keyboard, it just works
  \item MCP is the USB standard for AI tools
\end{itemize}
\end{shadedbox}
\end{column}
\end{columns}
\vspace{0.3cm}
You can connect multiple servers at once.  The LLM sees all their tools and decides which to use.
\end{frame}

\begin{frame}{MCP Server = Running Program}
An \textbf{MCP server} is a small program that stands between the AI and an external tool.
\vspace{0.3cm}
\begin{itemize}
  \item The server \alert{advertises tools} --- functions the AI can call (e.g., ``query stock prices'', ``get income statement'')
  \item When the AI decides to use a tool, the server \alert{makes the actual API call}, queries the database, or controls the browser
  \item The server returns structured results back to the AI
\end{itemize}
\vspace{0.3cm}
The AI never contacts the external service directly.  It sends a request to the MCP server; the server does the work and reports back.
\end{frame}

\begin{frame}{How MCP Works}
\begin{center}
\begin{tikzpicture}[
  node distance=1.2cm and 2.5cm,
  every node/.style={font=\small},
  sbox/.style={draw=titlegray, rounded corners, minimum width=3.2cm, minimum height=0.9cm, fill=excelinput, text=titlegray, font=\small\bfseries, align=center},
  sarrow/.style={-{Stealth[length=3mm]}, thick, draw=accentblue}
]
\node[sbox, minimum width=3.5cm] (app) {Claude Desktop\\(MCP Client)};
\node[sbox, right=2.5cm of app, yshift=1.5cm] (s1) {Browser\\Server};
\node[sbox, right=2.5cm of app] (s2) {Financial\\Data Server};
\node[sbox, right=2.5cm of app, yshift=-1.5cm] (s3) {Terminal\\Server};

\draw[sarrow] (app) -- (s1) node[midway, above, font=\scriptsize\itshape, text=titlegray, sloped] {MCP protocol};
\draw[sarrow] (app) -- (s2) node[midway, above, font=\scriptsize\itshape, sloped] {MCP protocol};
\draw[sarrow] (app) -- (s3) node[midway, below, font=\scriptsize\itshape, sloped] {MCP protocol};

\node[right=0.4cm of s1, font=\scriptsize, text=titlegray] {Chrome, Firefox, \ldots};
\node[right=0.4cm of s2, font=\scriptsize, text=titlegray] {FactSet, S\&P, Bloomberg, \ldots};
\node[right=0.4cm of s3, font=\scriptsize, text=titlegray] {PowerShell, Bash, \ldots};
\end{tikzpicture}
\end{center}
\vspace{0.3cm}
\begin{itemize}
  \item Claude Desktop is the \textbf{client}; each external service runs as a \textbf{server}
  \item Every server speaks the same protocol --- one standard, many tools
  \item Servers expose \textbf{tools} (functions the AI can call); the LLM picks which to use
\end{itemize}
\end{frame}

\begin{frame}{How an Agent Uses MCP}
\begin{center}
\begin{tikzpicture}[
  node distance=0.8cm and 2cm,
  every node/.style={font=\small},
  sbox/.style={draw=titlegray, rounded corners, minimum width=2.6cm, minimum height=1cm, fill=excelinput, text=titlegray, font=\small\bfseries, align=center},
  sarrow/.style={-{Stealth[length=3mm]}, thick, draw=accentblue}
]
\node[sbox] (user) {User};
\node[sbox, right=2cm of user, fill=accentblue!15] (agent) {Agent};
\node[sbox, right=2cm of agent] (llm) {LLM};
\node[sbox, below=1.2cm of agent, fill=alertorange!15] (mcp) {MCP\\Server};
\node[sbox, right=2cm of mcp] (tool) {External\\Tool};

\draw[sarrow] ([yshift=2pt]user.east) -- ([yshift=2pt]agent.west);
\draw[sarrow] ([yshift=-2pt]agent.west) -- ([yshift=-2pt]user.east);

\draw[sarrow] ([yshift=2pt]agent.east) -- ([yshift=2pt]llm.west);
\draw[sarrow] ([yshift=-2pt]llm.west) -- ([yshift=-2pt]agent.east);

\draw[sarrow] ([xshift=2pt]agent.south) -- ([xshift=2pt]mcp.north);
\draw[sarrow] ([xshift=-2pt]mcp.north) -- ([xshift=-2pt]agent.south);

\draw[sarrow] ([yshift=2pt]mcp.east) -- ([yshift=2pt]tool.west);
\draw[sarrow] ([yshift=-2pt]tool.west) -- ([yshift=-2pt]mcp.east);
\end{tikzpicture}
\end{center}
\vspace{0.2cm}
\begin{itemize}
  \item The \textbf{LLM} decides what to do and formats tool calls
  \item The \textbf{MCP server} executes the call --- the LLM never touches the API directly
  \item Results flow back through the agent to the user
\end{itemize}
\end{frame}

\begin{frame}{Why Not Just Call the API Directly?}
An agent \textit{can} access any API by writing code --- Claude Code does this all the time.  MCP servers are the alternative: a pre-built connector that standardizes the interface.
\vspace{0.3cm}
\begin{columns}[T]
\begin{column}{0.5\textwidth}
\begin{shadedbox}[title=\textbf{Direct API Access}]
\begin{itemize}\small
  \item Agent writes Python to call the API
  \item Must handle auth, parsing, errors each time
  \item Flexible but reinvents the wheel
\end{itemize}
\end{shadedbox}
\end{column}
\begin{column}{0.5\textwidth}
\begin{shadedbox}[title=\textbf{MCP Server}]
\begin{itemize}\small
  \item Pre-built connector, tested and shared
  \item Agent calls a tool; server handles the rest
  \item Works with any MCP-compatible client
\end{itemize}
\end{shadedbox}
\end{column}
\end{columns}
\vspace{0.3cm}
Like a USB driver: built once, works everywhere.  The agent doesn't need to know how the API works --- just what tools are available.
\end{frame}

\begin{frame}{What MCP Servers Handle for You}
\begin{itemize}
  \item \textbf{Authentication}: API keys, OAuth tokens, and credential refresh --- the LLM never sees secrets
  \item \textbf{Rate limiting}: manages API quotas and retries failed requests automatically
  \item \textbf{Data transformation}: cleans raw API responses into structured results the LLM can reason about
  \item \textbf{Security boundary}: controls exactly which operations are exposed (e.g., read-only database access)
  \item \textbf{Tool discovery}: advertises available tools with descriptions, so the LLM knows what it can do without being told in the prompt
\end{itemize}
\end{frame}

\begin{frame}{Three Ways to Install a Connection to an MCP Server}
\begin{columns}[T]
\begin{column}{0.33\textwidth}
\begin{shadedbox}[title=\textbf{Desktop Extensions}]
\begin{enumerate}\small
  \item Settings $\rightarrow$ Extensions
  \item Browse and find the server
  \item Click \textbf{Install}
\end{enumerate}
\vspace{0.2cm}
One-click. Anthropic-reviewed.
\end{shadedbox}
\end{column}
\begin{column}{0.33\textwidth}
\begin{shadedbox}[title=\textbf{Manual Config}]
\begin{enumerate}\small
  \item Install prerequisites (Node.js or uv)
  \item Edit JSON config file
  \item Restart Claude Desktop
\end{enumerate}
\vspace{0.2cm}
More control. Any server.
\end{shadedbox}
\end{column}
\begin{column}{0.33\textwidth}
\begin{shadedbox}[title=\textbf{Ask Claude Code}]
\begin{enumerate}\small
  \item Tell Claude Code which server you want
  \item It edits the config for you
  \item Restart Claude Desktop
\end{enumerate}
\vspace{0.2cm}
Easiest for Code users.
\end{shadedbox}
\end{column}
\end{columns}
\end{frame}

\begin{frame}[fragile]{You May Need to Install uv}
Some MCP servers are distributed as Python packages and run via \texttt{uvx}.  You may need to install \texttt{uv} first.
\vspace{0.3cm}
\begin{shadedbox}[title=\textbf{Manual install}]
\begin{itemize}\small
  \item \textbf{Mac}: \texttt{brew install uv}
  \item \textbf{Windows}: \texttt{winget install astral-sh.uv}
  \item Verify: \texttt{uv --version}
\end{itemize}
\vspace{0.1cm}
{\scriptsize Used by: Browser-Use, Alpha Vantage, and others}
\end{shadedbox}
\vspace{0.3cm}
If you use Claude Code to install a connection, it will ask permission to install \texttt{uv} automatically if needed.
\end{frame}

% ========================================
% SECTION 3: FINANCIAL DATA SERVERS
% ========================================
\section{Connecting to Financial Data}

\begin{frame}{Enterprise Financial Data via MCP}
Major financial data providers now offer official MCP servers, giving AI agents direct, governed access to institutional-grade data.
\vspace{0.3cm}
\begin{itemize}
  \item \textbf{FactSet} --- fundamentals, M\&A intelligence, real-time market data; first production-grade financial MCP server (800+ institutional users)
  \item \textbf{S\&P Capital IQ} --- financial statements, earnings transcripts, company fundamentals (built by Kensho, S\&P's AI subsidiary)
  \item \textbf{LSEG / Refinitiv} --- financial data via LSEG Workspace APIs
  \item \textbf{Bloomberg} --- community-built \texttt{blpapi-mcp} server (requires a Bloomberg Terminal); Bloomberg is an MCP governance member but has no official server yet
\end{itemize}
\end{frame}

\begin{frame}{Enterprise Financial Data via MCP (cont.)}
\begin{itemize}
  \item \textbf{Moody's} --- credit ratings, fundamentals, entity linkages, market indicators
  \item \textbf{Morningstar} --- ratings, fair value, economic moat, fund data, editorial research
  \item \textbf{PitchBook} --- private market data (PE, VC, M\&A)
  \item \textbf{Daloopa} --- financial metrics and KPIs for 5,000+ companies, every datapoint hyperlinked to its source
\end{itemize}
\vspace{0.3cm}
These are subscription services.  Install the server, ask Claude to pull data --- \alert{no code, no API documentation needed}.  The exercise uses free alternatives.
\end{frame}

\section{Browser Use}

\begin{frame}{Browser-Use: Web Automation via MCP}
\textbf{Browser-Use} is an MCP server that lets Claude autonomously control a web browser. Describe a task in plain English, and it handles the navigation.
\vspace{0.3cm}
\begin{columns}[T]
\begin{column}{0.5\textwidth}
\textbf{Finance Applications}
\begin{itemize}\small
  \item ``Search for AAPL on Yahoo Finance and get the P/E ratio''
  \item ``Download all linked CSV files from this SEC page''
  \item ``Fill out this application form with my information''
\end{itemize}
\end{column}
\begin{column}{0.5\textwidth}
\textbf{Key Features}
\begin{itemize}\small
  \item Plain-language task descriptions
  \item Navigates pages and fills forms
  \item Extracts data from web pages
  \item \alert{Recommended}: local (self-hosted) for full privacy
\end{itemize}
\end{column}
\end{columns}
\end{frame}

% ========================================
% SECTION 4: CONNECTING TO YOUR COMPUTER
% ========================================
\section{Terminal Control}

\begin{frame}{DesktopCommander: Terminal Control via MCP}
\textbf{DesktopCommander} gives Claude Desktop the ability to execute terminal commands, manage files, and control processes --- without switching to Code mode.
\vspace{0.3cm}
\begin{columns}[T]
\begin{column}{0.5\textwidth}
\begin{shadedbox}[title=\textbf{What It Can Do}]
\begin{itemize}\small
  \item Run terminal commands (\texttt{pip install}, git, scripts)
  \item Read, write, and search files
  \item Manage running processes
\end{itemize}
\end{shadedbox}
\end{column}
\begin{column}{0.5\textwidth}
\begin{shadedbox}[title=\textbf{When to Use It}]
\begin{itemize}\small
  \item You want to stay in Claude Desktop Chat
  \item Quick terminal tasks
  \item \alert{Always review commands before approving}
\end{itemize}
\end{shadedbox}
\end{column}
\end{columns}
\end{frame}

% ========================================
% SECTION 5: PERSONAL PRODUCTIVITY
% ========================================
\section{Personal Productivity}

\begin{frame}{Email and Calendar via MCP}
MCP servers can connect Claude to your personal productivity tools, turning it into an assistant that can \alert{read, search, and act} on your behalf.
\vspace{0.3cm}
\begin{columns}[T]
\begin{column}{0.5\textwidth}
\begin{shadedbox}[title=\textbf{Email (Gmail)}]
\begin{itemize}\small
  \item Search your inbox by sender, date, or keyword
  \item Draft and send replies
  \item Summarize unread messages
  \item ``Find all emails from our auditor about Q4 adjustments''
\end{itemize}
\end{shadedbox}
\end{column}
\begin{column}{0.5\textwidth}
\begin{shadedbox}[title=\textbf{Calendar (Google Calendar)}]
\begin{itemize}\small
  \item Check availability and schedule meetings
  \item List upcoming events
  \item Create, update, or cancel events
  \item ``What's on my calendar this week?''
\end{itemize}
\end{shadedbox}
\end{column}
\end{columns}
\end{frame}

\begin{frame}{Other Productivity Servers}
The MCP ecosystem includes servers for many tools you already use:
\vspace{0.3cm}
\begin{columns}[T]
\begin{column}{0.5\textwidth}
\begin{itemize}\small
  \item \textbf{Google Drive}: search and read documents, spreadsheets, and slides
  \item \textbf{Slack}: search channels, read messages, post updates
  \item \textbf{Notion}: query databases, read and update pages
\end{itemize}
\end{column}
\begin{column}{0.5\textwidth}
\begin{itemize}\small
  \item \textbf{GitHub}: manage repos, issues, and pull requests
  \item \textbf{Todoist / Linear}: manage tasks and projects
  \item \textbf{Memory}: give Claude persistent memory across conversations
\end{itemize}
\end{column}
\end{columns}
\vspace{0.3cm}
\alert{One prompt can span multiple tools}: ``Check my calendar for Friday, draft an email to the team with the agenda, and create a shared Google Doc for meeting notes.''
\end{frame}

% ========================================
% SECTION 6: COMBINING CONNECTIONS
% ========================================
\section{Combining Multiple Connections}

\begin{frame}[fragile,shrink=10]{Multi-Server MCP Configurations}
You can connect multiple MCP servers at once. Claude sees all their tools and can use them together.
\vspace{0.2cm}
\begin{verbatim}
{
  "mcpServers": {
    "browser-use": {
      "command": "uvx",
      "args": ["--from", "browser-use[cli]",
               "browser-use", "--mcp"],
      "env": {"ANTHROPIC_API_KEY": "sk-ant-..."}
    },
    "desktop-commander": {
      "command": "npx",
      "args": ["-y", "desktop-commander-mcp"]
    },
    "alphavantage": {
      "command": "uvx",
      "args": ["av-mcp", "YOUR_AV_KEY"]
    }
  }
}
\end{verbatim}
\end{frame}

\begin{frame}{Built-In Alternatives to MCP}
DesktopCommander and Browser-Use add capabilities to Claude's \textbf{Chat} mode. But Code and Cowork modes have some of these built in.
\vspace{0.3cm}
\begin{itemize}
  \item \textbf{Chat mode} needs MCP servers for terminal, file access, and browser control --- no internet access
  \item \textbf{Code tab} has terminal and file access built in, plus full internet --- browser control via Chrome extension
  \item \textbf{Cowork tab} has terminal and file access built in, but no internet --- browser control via Chrome extension
\end{itemize}
\vspace{0.3cm}
MCP servers for \textbf{financial data} (Alpha Vantage, FMP) and \textbf{personal productivity} (email, calendar, Google Drive) remain valuable in all modes.
\end{frame}

% ========================================
% SECTION 6: PLUGINS AND THE SAASPOCALYPSE
% ========================================
\section{Plugins and the SaaSpocalypse}

\begin{frame}{What is a Plugin?}
A \textbf{plugin} is a packaged collection of skills, agents, hooks, and MCP servers that can be shared across teams.  (\textit{Hooks} are shell commands that run automatically before or after specific events --- e.g., running a linter every time Claude edits a file.)
\vspace{0.3cm}
\begin{columns}[T]
\begin{column}{0.5\textwidth}
\textbf{A Plugin Can Include}
\begin{itemize}\small
  \item Skills (\texttt{SKILL.md} files)
  \item Agents (\texttt{AGENT.md} files)
  \item MCP servers (external tools)
\end{itemize}
\end{column}
\begin{column}{0.5\textwidth}
\textbf{Why Use Plugins?}
\begin{itemize}\small
  \item Share skills with your team
  \item Namespace prevents conflicts
  \item Install from marketplaces
\end{itemize}
\end{column}
\end{columns}
\vspace{0.3cm}
\textbf{Practical workflow:} Start with a standalone skill for quick prototyping. When ready to share, wrap it in a plugin by adding a manifest file.
\end{frame}

\begin{frame}[shrink=10]{The ``SaaSpocalypse'': Market Reaction to Plugins}
On January 30, 2026, Anthropic released 11 open-source plugins for Claude.  Within days, \alert{\$285 billion} in market cap evaporated from software, legal tech, and financial services stocks.
\vspace{0.2cm}
\begin{columns}[T]
\begin{column}{0.5\textwidth}
\textbf{Hardest-Hit Stocks}
\begin{itemize}\small
  \item LegalZoom: $-$20\%
  \item Thomson Reuters: $-$16\%
  \item RELX (LexisNexis): $-$14\%
  \item Intuit: $>-$10\%; Salesforce: $-$7\%
\end{itemize}
\end{column}
\begin{column}{0.5\textwidth}
\textbf{Why It Happened}
\begin{itemize}\small
  \item Plugins showed AI could replicate \textit{core workflows} of specialized enterprise software
  \item \textbf{No proprietary code} --- just markdown instructions + database connectors
  \item Investors asked: ``Why pay \$50K/yr for software when a plugin does it?''
\end{itemize}
\end{column}
\end{columns}
\end{frame}

% ========================================
% SECTION 7: EXERCISES
% ========================================
\section{Exercises}

\begin{frame}[fragile,shrink=5]{Example: Install a Financial Data Server}
Free MCP servers let you try this workflow without a subscription:
\vspace{0.2cm}
\begin{columns}[T]
\begin{column}{0.5\textwidth}
\begin{shadedbox}[title=\textbf{Alpha Vantage}]
\begin{itemize}\scriptsize
  \item Stock prices, options, forex, crypto, economic indicators
  \item Free key at \url{https://www.alphavantage.co} (25 req/day)
\end{itemize}
\vspace{0.1cm}
\begin{verbatim}
"alphavantage": {
  "command": "uvx",
  "args": ["av-mcp", "YOUR_KEY"]
}
\end{verbatim}
\end{shadedbox}
\end{column}
\begin{column}{0.5\textwidth}
\begin{shadedbox}[title=\textbf{Financial Modeling Prep}]
\begin{itemize}\scriptsize
  \item Financial statements, fundamentals, SEC filings
  \item Free key at \url{https://financialmodelingprep.com} (250 req/day)
\end{itemize}
\vspace{0.1cm}
\begin{verbatim}
"financial-modeling-prep": {
  "command": "npx",
  "args": ["-y", "@houtini/fmp-mcp"],
  "env": {"FMP_API_KEY": "YOUR_KEY"}
}
\end{verbatim}
\end{shadedbox}
\end{column}
\end{columns}
\vspace{0.2cm}
\begin{enumerate}\small
  \item Get a free API key and install the connection in Claude Desktop
  \item Ask Claude to retrieve financial data, analyze trends, and produce formatted Excel output
\end{enumerate}
\end{frame}

\begin{frame}{Exercise: Browser Automation}
Use Browser-Use MCP to navigate a financial website:
\begin{itemize}
  \item Extract market cap, P/E ratio, and revenue for a company
  \item Compile a summary comparison of 3 companies
  \item Save results to a formatted document
\end{itemize}
\end{frame}

\begin{frame}{Summary}
\begin{columns}[T]
\begin{column}{0.33\textwidth}
\begin{shadedbox}[title=\textbf{MCP}]
\begin{itemize}\small
  \item Universal plug-in protocol
  \item Financial data, browser, terminal
  \item Works with Claude Desktop
\end{itemize}
\end{shadedbox}
\end{column}
\begin{column}{0.33\textwidth}
\begin{shadedbox}[title=\textbf{Code Execution}]
\begin{itemize}\small
  \item Chat, Cowork, Code modes
  \item Match the tool to the task
  \item Code for APIs and files
\end{itemize}
\end{shadedbox}
\end{column}
\begin{column}{0.33\textwidth}
\begin{shadedbox}[title=\textbf{AI in Apps}]
\begin{itemize}\small
  \item Excel, Google Sheets
  \item Formulas, not hardcodes
  \item Choose by ecosystem
\end{itemize}
\end{shadedbox}
\end{column}
\end{columns}
\vspace{0.5cm}
\begin{center}
\alert{AI becomes truly powerful when it can reach your data, your tools, and your files.}
\end{center}
\end{frame}

\end{document}
